\documentclass[11pt]{article}
\usepackage[letterpaper,margin=1in]{geometry}
\usepackage{amsmath,amssymb,amsthm,booktabs,microtype}
\usepackage[hidelinks]{hyperref}
\newtheorem{theorem}{Theorem}
\DeclareMathOperator{\supp}{supp}
\newcommand{\dd}{\mathrm d}
\newcommand{\ii}{\mathrm i}
\newcommand{\e}{\mathrm e}
\title{The sixth-order cost of splitting a zero-cancelling filter}
\author{}
\date{}
\begin{document}
\maketitle
\begin{abstract}
For positive averages cancelling prescribed frequencies, we prove a sharp
lower bound for the support diameter in terms of their arithmetic mean.
A strip identity characterizes equality and measures the excess as an
integrated loss of phase coherence. Near a multiple frequency this excess
is comparable to the sum of the squares of the odd elementary symmetric
functions of the centered frequencies. Every asymmetric splitting direction
has an exact even order of vanishing. The first term in the full expansion
has degree six and an explicit coefficient. These conclusions answer the
ProveIt article's question about the cost and sign of asymmetric splitting.
\end{abstract}

\section{The extremal problem}
For a compactly supported probability measure $\mu$ on $\mathbb R$ write
$F_\mu(\lambda)=\int\exp(\ii\lambda t)\,\dd\mu(t)$.
Given a multiset of positive frequencies $\lambda_1,\dots,\lambda_N$, let
$B(\boldsymbol\lambda)$ be the minimum diameter of $\supp\mu$ when
$F_\mu$ vanishes at the specified frequencies, to the combined order of
each repetition. Define
\[
 \kappa=N^{-1}\sum_j\lambda_j,\qquad
 u_j=\lambda_j/\kappa-1,\qquad d=\max_j |u_j|,
 \qquad e_j=e_j(u_1,\dots,u_N)
\]
with $e_j$ the elementary symmetric functions. In particular,
$e_1=0$ and $e_3=\tfrac13\sum_j u_j^3$.

\begin{theorem}[Global bound and equality]\label{global}
For every positive frequency multiset, $\kappa B\ge N\pi$.
Set $h=\pi/\kappa$. Equality holds precisely when all
$\lambda_j<2\kappa$, the multiset is invariant under
$\lambda\mapsto2\kappa-\lambda$, and the coefficients of
$A(z)=\prod_j(z-\exp(\ii h\lambda_j))$ are nonnegative.
In this case the coefficients of $A(z)/A(1)$ are the masses of an optimizer
on $0,h,\ldots,Nh$.
\end{theorem}

\begin{theorem}[Local rigidity and the first correction]\label{main}
Fix $N\ge3$ and put $S=\sum_{3\le j\le N,\;j\ {\rm odd}} e_j^2$.
There are constants $0<c_N<C_N$ such that, for all sufficiently small $d$,
\[
 c_NS\le \kappa B(\boldsymbol\lambda)-N\pi\le C_NS .
\]
On each fixed centered asymmetric ray $u_j=tv_j$ the excess is asymptotic
to a strictly positive multiple of $t^{2r}$, where $r$ is the least odd index
with $e_r(v)\ne0$. For every cluster the uniform Taylor expansion starts with
\[
 \boxed{\displaystyle
 \kappa B(\boldsymbol\lambda)=N\pi+
 \frac{\pi^5 N}{20(N-1)(N-2)}\,e_3^2+O_N(d^8).}
\]
\end{theorem}

The repository proves $\kappa B=N\pi+O_N(d^3)$ and exact equality for
sufficiently tight reflection-invariant clusters. It asks for the first
nonzero correction and its sign on nonsymmetric clusters [1,
\emph{Splitting a multiple frequency} and \emph{Open questions}].
Theorem~\ref{main} identifies that correction even along directions on which
the degree-six term vanishes. In ranks three and four, $e_3$ is the only
possible nonzero odd coefficient. In ranks one and two every centered multiset
is reflection-invariant and the local formula is exactly $\kappa B=N\pi$.

For a zero $\rho=\beta+\ii\gamma$ in an explicit formula, evaluating a prime
sum at $x\e^t$ multiplies its $x^\rho$ term by $\exp(\rho t)$.
Positive tilting $\dd\nu=\e^{\beta t}\dd\mu/\int\e^{\beta v}\dd\mu(v)$ identifies
the multiplier and all its zero conditions with the Fourier conditions for
$\nu$ at $\gamma$. The two measures have identical support and the tilt is
invertible. Thus Theorem~\ref{main} gives the exact leading cost of cancelling
a tight cluster on any one vertical line, including Dirichlet zero ordinates
under the repository's GRH input. The extremal theorem itself uses no
arithmetic hypothesis. Modal input shaping also studies the minimum duration
of positive impulse sequences that cancel prescribed modes [2].


In absolute-height units, put $H=\max_j|\lambda_j-\kappa|$. For $H/\kappa$
sufficiently small the minimum log-width is
\[
 B=\frac{N\pi}{\kappa}
   +\frac{\pi^5N}{180(N-1)(N-2)}
       \frac{\bigl(\sum_j(\lambda_j-\kappa)^3\bigr)^2}{\kappa^7}
   +O_N\!\left(\frac{H^8}{\kappa^9}\right).
\]
In particular, any $N$ prescribed zeta zero ordinates in a window of fixed
length at height $T$ have a filter of optimal width
$N\pi/\kappa+O_N(H^6/\kappa^7)$ as $T$ tends to infinity. Such windows
occur at every sufficiently large height if their fixed length is chosen
large enough: the Riemann--von Mangoldt formula
$N_\zeta(t)=\tfrac{t}{2\pi}\log(t/(2\pi e))+O(\log t)$
gives at least $N$ ordinates, counted with multiplicity, in $[T,T+C]$
by taking the constant $C$ large enough to dominate the two remainder
terms. Repeated selected ordinates specify repeated vanishing. On the
critical line the corresponding exponentially tilted weights cancel
these modes in the prime explicit formula.

\section{A strip identity and the equality case}
For a finite probability distribution with positive endpoint weights on
$0=t_0<\cdots<t_M=L$, set $G(s)=\sum w_j\exp(t_js)$.
The triangle inequality gives $J(T)\ge0$ for
\begin{equation}\label{strip}
 J(T)=\int_{\mathbb R}\log\frac{G(x)}{|G(x+\ii T)|}\,\dd x
 =\frac{LT^2}{2}-2\pi\sum_{0<\Im\rho<T}(T-\Im\rho).
\end{equation}
The sum ranges over the zeros of $G$ with their orders. To verify the
identity, each end of the exponential sum has a single dominant term
uniformly in $\Im s$, so the integrand decays exponentially in $|x|$ and
zeros have bounded real part. Logarithmic singularities are integrable.
In distributions, $\Delta\log|G|=2\pi\sum_\rho\delta_\rho$, by local
factorization and $\Delta\log|s|=2\pi\delta_0$. Integrating the $xx$
derivative in $x$ gives the difference of the endpoint slopes $L-0$,
so $J''(T)=L-2\pi\sum_\rho\delta_{\Im\rho}(T)$.
Integration on $[0,T]$ gives \eqref{strip}, since $J(0)=0$ and
$J'(0)=0$ by conjugation and zero-freeness near the real axis.
One can first truncate the $x$ integral outside the zero-containing
strip and then take the limit using the exponential tails.
This is the exponential-sum analogue of the integrated angular inequality
of Eremenko and Fryntov for nonnegative-coefficient polynomials [3].

An optimum exists because a convolution of two-point filters is feasible and
an anchored minimizing sequence has a weakly convergent subsequence on a
compact interval. For any feasible measure the real and imaginary jet vector
is continuous on its support. Its mean lies in the convex hull of its
values there, so Carath\'eodory's theorem replaces it by a finite measure
with the same constraints and no larger diameter. Translate such an
optimal finite measure to start at zero. At $T=2\kappa$, the zero sum in
\eqref{strip} is at least
$\sum_j(T-\lambda_j)_+\ge NT-\sum_j\lambda_j=N\kappa$.
Nonnegativity gives $L\ge N\pi/\kappa$.
Equality forces $J(T)=0$ and $\lambda_j\le T$ for every $j$.
At every real $x$ the triangle inequality is then an equality, by
continuity (a zero on this line would instead give a positive integral).
At $x=0$ the phases are therefore all one since $t_0=0$, so every node is in
$h\mathbb Z$ and $G(\ii T)=1$. This excludes $\lambda_j=T$. Now $G(s)$ is
a real polynomial in $\exp(hs)$ of degree exactly $N$.
The exponential map is injective for the $N$ targets in the open strip
$0<\Im s<T$ on the imaginary axis and is locally invertible, so its
polynomial must be $A/A(1)$. Real coefficients give reflection
invariance by conjugation; positivity of the leading coefficient and
nonnegative masses give nonnegative coefficients for $A$.
Conversely those coefficients realize the stated width. This proves
Theorem~\ref{global}.

\section{Why all terms below degree six disappear}
Rescaling time reduces the proof to $\kappa=1$. To avoid differentiating
individual roots at a multiple root, use the coefficients of
\[
 P(x)=\prod_j(x-u_j)=x^N+c_2 x^{N-2}+\cdots+c_N,
 \qquad c_j=(-1)^j e_j .
\]
For each analytic $f$ near the origin, its degree-$<N$ interpolation remainder
modulo $P$ depends analytically on the $c_j$. Indeed it is
\[
 R_P f(x)=\frac1{2\pi\ii}\int_C
 \frac{f(v)}{P(v)}\frac{P(v)-P(x)}{v-x}\,\dd v,
\]
where a fixed small circle $C$ encloses all roots. Cauchy's formula shows
that $f-R_P f$ is divisible by $P$, also for repeated roots.

The repository's cluster theorem [1] gives an optimal measure
$\sum_{j=0}^N w_j\delta_{t_j}$ with $t_0=0$, positive weights and
$t_N=B$. At coalescence its nodes are $j\pi$ and its weights are
$w_j^0=2^{-N}\binom Nj$. The $2N+1$ equations for its $2N+1$ free
weights and nodes are the mass equation and the real and imaginary parts of
the coefficients of $R_P(F_\mu(1+\cdot))=0$.
At $P=x^N$ these coefficients are precisely the first $N$ Taylor jets.
To see invertibility directly, put $z=\exp(\ii\pi x)$, $y=z-1$ and
$\ell=\log z$. At the base a real variation of the transform is
$A(z)+(\ii+\ell/\pi)T(z)$ with $A,T$ real polynomials of degree at most
$N$, $A(-1)=0$ from the mass condition and $T(0)=0$ from anchoring.
Its vanishing modulo $y^N$ forces $T=0$ by taking the imaginary part
as a formal series in real $y$, and then $A=0$ by taking the real part.
The node weights are all positive. The analytic implicit function theorem
therefore gives a solution in the coefficient variables on a full real
neighborhood, even at polynomials whose roots cease to be real.
At a nearby multiset of real roots this solution is the optimizer of [1]
by local uniqueness and continuity. Denote its last node by $b(c)$.

Write the coefficient vector as $(E,O)$ with its even and odd indices
grouped. If $O=0$, real simple mirror configurations form a relatively open
subset of the even coefficient space, arbitrarily near the origin, by taking
distinct pairs $\pm a_j$ and also $0$ when $N$ is odd. On this set
$b=N\pi$ by [1] because all weights remain strictly positive. The real
analytic identity theorem extends this identity to all nearby even
coefficient vectors.

Take $G_c(s)=\sum_jw_j(c)\exp(t_j(c)s)$ for any nearby real coefficient vector.
At $c=0$ the only zeros with $0<\Im s<2$ are the $N$ copies of $\ii$.
For nearby coefficients the endpoint dominance on $|\Re s|>R$ is uniform
because the endpoint weights exceed a fixed $w>0$ and adjacent gaps exceed a
fixed $g>0$. For instance, at $x>R$ division by the last term bounds the
sum of the remaining absolute values by $w^{-1}\exp(-gx)$; the first term
gives the analogous bound at $x<-R$.
Rouch\'e's theorem on the remaining compact rectangle gives precisely $N$
zeros there, the prescribed points $\rho_j=\ii(1+u_j)$ with multiplicity.
Here $u_j$ may be nonreal roots of $P$, and still
$\sum\Im\rho_j=N+\Re\sum u_j=N$. Formula \eqref{strip} therefore gives
\begin{equation}\label{defect}
 b(c)-N\pi=\tfrac12\int_{\mathbb R}
   -\log\left|\sum_j p_j(x,c)\exp(2\ii t_j(c))\right|\,\dd x,\qquad
 p_j(x,c)=\frac{w_j(c)\exp(xt_j(c))}{G_c(x)} .
\end{equation}
In particular $b\ge N\pi$ on a full real coefficient neighborhood, and
$\partial_O b(E,0)=0$. Its Taylor series relative to $N\pi$ has at least
two odd coefficients in every monomial.

\section{A positive quadratic form for asymmetric perturbations}
Consider $c=\epsilon(0,v)$ with $v$ supported on odd indices, and put
$s_j=\partial_\epsilon t_j|_0$. The coefficient of $\epsilon^2$ in
the right side of \eqref{defect} is
\begin{equation}\label{quad}
 Q(v)=\int_{\mathbb R}\operatorname{Var}_{p(x,0)}(s_j)\,\dd x.
\end{equation}
Indeed the phases are all one at $\epsilon=0$. For any probability vector
$p$, the expansion $-\log|\sum p_j e^{2\ii h_j}|
=2\operatorname{Var}_p(h)+O(\max|h_j|^3)$ follows by Taylor expansion.
Derivatives of $p$ do not affect the quadratic coefficient since its mass
is one. Differentiation under the integral is valid: on each tail one
endpoint has all but $O(\exp(-g|x|))$ of the mass, uniformly for small
$c$, for a fixed positive $g$. Factoring out this endpoint phase gives
an integrand with two derivatives bounded there by
$C(1+|x|)^2\exp(-g|x|)$. The central compact interval is zero-free at
height two for all nearby $c$.

The form $Q$ is positive definite in the odd coordinates.
Indeed $Q(v)=0$ requires all $s_j$ equal and consequently zero since
$s_0=0$. At the base $F_0=Cy^N$, $C=(-1)^N2^{-N}$. Differentiating
$F_\mu(1+x)=P(x)Q_c(x)$ now shows that the first variation satisfies
\begin{equation}\label{firstvar}
 F_1\equiv C y^N\sum_{r\ {\rm odd}} v_r/x^r\pmod{x^N}.
\end{equation}
If all nodes have zero variation, $F_1$ is a polynomial in $z$ with real
coefficients. As a formal series in real $y$, the right side of
\eqref{firstvar} is purely imaginary, since $x=\ell/(\ii\pi)$ and $r$ is
odd. Its lowest degree is $N-r_*$ with a nonzero coefficient when $r_*$
is the largest index with $v_{r_*}\ne0$. This contradiction proves the
claim.

The Hessian in $O$ is consequently positive definite at zero and remains
uniformly positive definite in a smaller coefficient neighborhood.
Taylor's formula along the segment $(E,tO)$, using $b(E,0)=N\pi$ and
$\partial_O b(E,0)=0$, proves the two-sided bound in Theorem~\ref{main}.
It also gives $b-N\pi=Q(O)+o(\|O\|^2)$ as $(E,O)\to0$ with $O\ne0$.
On a fixed asymmetric ray the leading constant is
$Q(\mathbf e_r)e_r(v)^2$, where $\mathbf e_r$ is the coefficient coordinate
vector, proving the exact order. Giving $c_j$ weight $j$, the only possible
term of degree six is $c_3^2$ and no term has degree seven. All higher
terms are $O_N(d^8)$ by analyticity and $c_j=O_N(d^j)$.

\section{Evaluation of the sixth-order coefficient}
Set $v_3=1$, $v_j=0$ otherwise, and $b_j=(-1)^jw_j^0$. In the notation
above, $F_1=W(z)+(\ii+\ell/\pi)T(z)$ with $W$ real and
$T(z)=\sum_j b_j s_jz^j$. The imaginary part of \eqref{firstvar}
gives $T\equiv-\pi^3 C y^N/\ell^3\pmod{y^N}$.
Since $(y/\ell)^3=1+3y/2+y^2/2+O(y^3)$ and $T(0)=0$ at $z=0$,
the unique degree-at-most-$N$ solution is
\[
 T(z)=-\tfrac12\pi^3 C y^{N-3}z(z+1),\qquad
 f(j):=\frac{s_j}{\pi^3}
   =\frac{j(N-j)(2j-N)}{2N(N-1)(N-2)} .
\]
The last equality follows by extracting binomial coefficients.
Under $p(x,0)$ the node index has law $K\sim{\rm Bin}(N,p)$ with
$p=\exp(\pi x)/(1+\exp(\pi x))$. Thus
\[
 Q(\mathbf e_3)/\pi^5=\int_0^1\frac{\operatorname{Var}(f(K))}{p(1-p)}\,\dd p .
\]
For $m(p)=\mathbb E f(K)$ one has
\begin{equation}\label{bin}
 \operatorname{Var}(f(K))=
 \sum_{r=1}^N\frac{[p(1-p)]^r}{r!(N)_r}\,(m^{(r)}(p))^2,
 \qquad (N)_r=N(N-1)\cdots(N-r+1).
\end{equation}
Indeed write $K=\sum X_i$ for independent Bernoulli variables.
The multilinear interpolation of $f(\sum X_i)$ about the vector
$(p,\dots,p)$ is an exact sum in the orthogonal products
$\prod_{i\in S}(X_i-p)$. Its common coefficient for $|S|=r$ is
$m^{(r)}(p)/(N)_r$, by differentiating the symmetric multilinear
interpolant on the diagonal. Summing the squared norms gives \eqref{bin}.

The binomial factorial moments give $m(p)=p(1-p)(2p-1)/2$.
Only the first three terms of \eqref{bin} survive. Integrating
$[p(1-p)]^{r-1}(m^{(r)}(p))^2/r!$ gives $1/20,3/20,1/5$ respectively,
and hence
\[
 Q(\mathbf e_3)/\pi^5=
 \frac{1}{20N}+\frac{3}{20N(N-1)}
 +\frac{1}{5N(N-1)(N-2)}
 =\frac{N}{20(N-1)(N-2)}.
\]
This finishes the proof of Theorem~\ref{main}.

\section{A numerical example and the remaining question}
For $\boldsymbol\lambda=(1-2t,1-t,1+3t)$ we have $e_3=6t^3$.
The predicted limit of $(B-3\pi)/(36t^6)$ is $3\pi^5/40\approx22.95148$.
Independently solve the three frequency equations
$p\cos(v\lambda_j)+(1-p)\cos(w\lambda_j)=0$ near
$(v,w,p)=(3\pi/2,\pi/2,1/4)$ using symmetric nodes $\pm v,\pm w$.
The diameter is $2v$. Here the last column divides the observed coefficient
by the predicted limit.
\begin{center}
\begin{tabular}{rrr}\toprule
$t$&$(2v-3\pi)/(36t^6)$&ratio\\\midrule
$1/20$&24.35608839&1.06119920\\
$1/40$&23.30044786&1.01520475\\
$1/80$&23.03910421&1.00381796\\
$1/160$&22.97347516&1.00095849\\
$1/320$&22.95699039&1.00024025\\\bottomrule
\end{tabular}
\end{center}
The accompanying \texttt{check\_filter.py} verifies the rational identity
symbolically for arbitrary $N$ and reproduces the example at 90-decimal
working precision. The local analytic theorem supplies optimality for
sufficiently small $t$. The tabulated solutions illustrate convergence;
their particular radii of global optimality have not been certified.

An exceptional ray in rank five is $u=tv$ with
$v=(-9,-5,-1,7,8)$. Its entries and their cubes sum to zero, and
$e_5(v)=-2520$. Its first correction has degree ten:
\[
 \kappa B-5\pi\sim \frac{25\pi^9}{6912}\,2520^2 t^{10}.
\]
For a unit perturbation in $c_5$, \eqref{firstvar} gives
$T/(\pi^5 C)=z(z+1)(z^2+10z+1)/24$ by Taylor truncation and anchoring.
The node variations divided by $\pi^5$ are
$(0,1/120,-11/240,11/240,-1/120,0)$.
Their binomial mean equals $p(p-1)(2p-1)(12p^2-12p+1)/24$;
substitution in \eqref{bin} gives the stated coefficient.

The variance integral also gives the coefficients on the other exceptional rays.
A remaining problem is to evaluate $Q(\mathbf e_r)$ and its mixed entries
in a closed form uniform in $r$ and $N$. The global equality criterion
reduces the exact boundary of the equality region within the mirror locus
to coefficient nonnegativity of an explicit polynomial.

\begin{thebibliography}{9}
\bibitem{repo} V. Reshetnikov, \emph{ProveIt, Generalized Riemann Hypothesis
article}, section \emph{Positive filters that cancel zeros}. See the results
\emph{Reflection-paired spectra}, \emph{Differential at a mirror configuration}
and \emph{Splitting a multiple frequency}, and the second item of the section's
open questions.
\href{https://github.com/VladimirReshetnikov/ProveIt/blob/00cadb961bbe77c74e5d082c4948f8884d58a970/Analysis/GeneralizedRiemannHypothesis/Article/sections/fil.tex}{Pinned source on GitHub}.
\bibitem{pao} M. A. Lau and L. Y. Pao, \emph{Input shaping and time-optimal
control of flexible structures}, Automatica 39 (2003), 893--900.
\href{https://doi.org/10.1016/S0005-1098(03)00024-4}{Publisher record}.
\bibitem{ef} A. Eremenko and A. Fryntov, \emph{Remarks on the Obrechkoff
inequality}, Theorem 1, \href{https://arxiv.org/abs/1409.4849}{arXiv:1409.4849}.
\end{thebibliography}
\end{document}
